\documentclass[runningheads]{llncs}
\usepackage[T1]{fontenc}
\usepackage{graphicx}
\usepackage{booktabs}
\usepackage{amsmath}
\usepackage{subcaption}
\usepackage{hyperref}

\begin{document}

\title{Sneakie: Learning to Play the Snake Game in All Its Forms}
\titlerunning{Sneakie: Deep RL on Snake}

\author{JAMAL Adonis \and  KAPOTOS Fotios \and MARTINI Jean }
% \authorrunning{Kapotos et al.}
\institute{CentraleSupélec}
\maketitle

% ── Abstract ──────────────────────────────────────────────────────────────────
\begin{abstract}
We present a comparative study of deep reinforcement learning methods on Snake,
a grid-world game that exhibits many of the challenges found in real sequential
decision-making problems: sparse rewards, partial observability, long-horizon
planning, and increasingly complex environment dynamics.
We develop a Snake environment that supports multiple food types (gold, silver, and poison), static and dynamic obstacles,and three distinct observation spaces: a 15-dimensional hand-crafted feature
vector, an egocentric convolutional window, and a full-grid CNN observation.
Using DQN and its extensions (Double DQN, Dueling networks, frame stacking), we study the ability of our agents to solve, generalize and adapt to a range of different environments, from easy to complex.
Our results show that CNN-based agents consistently outperform hand-crafted
MLP agents in complex environments.
\keywords{Deep Q-Network \and Reward Shaping \and State Representation
          \and Snake \and Gymnasium}
\end{abstract}

% ── 1. Introduction ───────────────────────────────────────────────────────────
\section{Introduction}
\label{sec:intro}

The game of Snake is a simple environment that captures key challenges in
reinforcement learning: the agent must collect food to grow while avoiding
collisions with walls, obstacles, and its own body.
To increase the richness of the environment, we introduce a \emph{hard mode}
with multiple food types (gold, silver, poison) and dynamic obstacles.

We study how different state representations and training strategies impact
performance across environments of increasing complexity.

We address three research questions:
\begin{enumerate}
  \item \textbf{Solving easy mode.} How do hand-crafted features compare to
        learned spatial representations (egocentric CNN vs.\ full-grid CNN)?
  \item \textbf{Reward design.} What is the effect of dense reward signals
        (distance shaping, body-repulsion) on convergence and performance?
  \item \textbf{Scalability to hard mode.} How do different architectures
        generalise as environment complexity increases?
\end{enumerate}

\noindent\textbf{Contributions.}
\begin{itemize}
  \item A configurable Snake environment with multiple observation spaces,
        reward functions, and a hard mode with dynamic obstacles.
  \item A comparison of hand-crafted features and CNN-based representations.
  \item An evaluation of DQN extensions (Double DQN, Dueling networks,
        frame stacking) in complex environments.
\end{itemize}

% ── 2. Environment ────────────────────────────────────────────────────────────
\section{Environment}
\label{sec:env}

\subsection{State Representation}
\label{sec:obs}

The environment is a discrete $H \times W$ grid (typically $16 \times 16$).
Each cell encodes one of the following values: empty cell, snake body, snake head, gold food, silver food, poisson or obstacle.
We consider three observation modalities:

\textbf{Hand-crafted features (MLP).}
A 15-dimensional binary vector encoding:
(i) collision risks (straight, left, right),
(ii) current direction (one-hot),
(iii) relative position of nearest food,
(iv) relative position of nearest poison.

\textbf{Egocentric CNN.}
A $C \times k \times k$ tensor centred on the head, with $C=6$ channels
(body, head, gold, silver, poison, obstacles).
This enforces locality and translation invariance.

\textbf{Full-grid CNN.}
A $C \times H \times W$ tensor covering the entire board,
providing global spatial information at higher dimensional cost.

\subsection{Action Space and Dynamics}
\label{sec:action}

At each step, the agent selects one of four \emph{absolute} actions:
up, right, down, or left.

Reverse moves can optionally be disallowed, in which case attempting
to reverse direction results in continuing forward.

The snake evolves as follows:
\begin{itemize}
  \item The head moves to a new cell according to the chosen action.
  \item If no food is consumed, the tail is removed (constant length).
  \item If gold or silver food is consumed, the snake grows (tail kept).
  \item If poison is consumed, the snake shrinks by a fixed number of segments.
\end{itemize}
\subsection{Reward Function}
\label{sec:reward}

Rewards are additive:
\begin{itemize}
  \item Gold: $+r_g$, \quad Silver: $+r_s$
  \item Poison: $+r_p$ ($\leq 0$), with shrink
  \item Death: $r_d < 0$, \quad Step: $r_{\text{step}}$
\end{itemize}

When no food is consumed, we add potential-based shaping:
\begin{align}
r_t \mathrel{+}= \lambda (d_{t-1} - d_t) + \mu (b_t - b_{t-1}),
\end{align}
where $d_t$ is the Manhattan distance to the nearest (non-poison) food and
$b_t$ the distance to the nearest body segment.

The first term drives food-seeking, while the second discourages
self-collisions, especially for long snakes.

\subsection{Obstacles and Dynamics}
\label{sec:obstacles}

The environment includes three types of obstacles:

\begin{itemize}
  \item \textbf{Static obstacles}: fixed cells placed at initialization.
  \item \textbf{Random obstacles}: sampled at each episode reset.
  \item \textbf{Dynamic obstacles}: 3-cell walls that oscillate deterministically
        along a fixed axis (horizontal or vertical), reversing direction every step.
\end{itemize}

Food is never spawned in cells occupied by obstacles or in regions swept by
dynamic obstacles.

\subsection{Environment Configurations}
\label{sec:configs}

We define two main environment regimes used throughout our experiments.

\textbf{Easy environment.}
\begin{itemize}
  \item Grid size: $16 \times 16$
  \item Food: 1 gold item (no silver or poison)
  \item Obstacles: none
  \item Dynamics: static environment
\end{itemize}
This configuration isolates core navigation and growth behaviour without additional sources of stochasticity or risk.

\textbf{Hard environment.}
\begin{itemize}
  \item Grid size: $16 \times 16$
  \item Food: multiple items (typically 2 gold, 3 silver, 1 poison)
  \item Obstacles:
  \begin{itemize}
    \item Static obstacles (fixed walls)
    \item Random obstacles (sampled at reset)
    \item Dynamic obstacles (3-cell walls oscillating along a fixed axis)
  \end{itemize}
  \item Dynamics: non-stationary due to moving obstacles
\end{itemize}

This setting introduces long-horizon planning challenges, partial observability under egocentric views, and increased risk of self-trapping, particularly for long snakes.
% ── 3. Methods ────────────────────────────────────────────────────────────────
\section{Methods}
\label{sec:methods}

All agents share the same DQN training loop: $\epsilon$-greedy exploration with
linear annealing, experience replay (buffer size 100\,000), and a target
network updated every $C$ steps.

\subsection{MLP DQN (Baseline)}
\label{sec:mlp}

A 3-layer fully-connected network ($256$--$128$--$3$) consumes the 15-dim
hand-crafted feature vector.
This is the simplest and most data-efficient architecture — the features
encode exactly the information needed for collision avoidance and food targeting
on the simple environment.

\subsection{CNN DQN — Egocentric Window}
\label{sec:cnn_window}

Two convolutional layers followed by two fully-connected layers process the
$6 \times k \times k$ egocentric observation.
The local receptive field and translation equivariance match the structure of
the control problem: the decision of whether to turn depends primarily on
what is immediately around the head.

\subsection{CNN DQN — Full Grid}
\label{sec:cnn_full}

The same convolutional architecture is applied to the full $C \times H \times W$
board observation.
The agent retains access to global context (distant food, global body shape),
which is expected to help in long-horizon planning at the cost of a higher-
dimensional input.

\subsection{Double DQN and Dueling Networks}
\label{sec:ddqn_dueling}

\textbf{Double DQN} decouples action selection from action evaluation,
reducing the overestimation bias of vanilla DQN:
\begin{equation}
  y_t = r_t + \gamma\, Q\bigl(s_{t+1},\,\arg\max_a Q(s_{t+1}, a;\,\theta);\,\theta^-\bigr).
\end{equation}

\textbf{Dueling networks} split the convolutional trunk into a state-value head
$V(s)$ and an advantage head $A(s,a)$, combining them as
$Q(s,a) = V(s) + \bigl(A(s,a) - \overline{A}(s)\bigr)$.
This decomposition accelerates convergence in states where survival matters
more than which food to pursue next.

\subsection{Frame Stacking ($n=4$)}
\label{sec:framestack}

The last four grid observations are concatenated into a
$(4C) \times H \times W$ tensor before being fed to the CNN.
Frame stacking encodes obstacle velocity from positional differences across
frames, converting the partially-observable dynamic-obstacle setting into an
approximately Markovian state.

% ── 4. Experiments ────────────────────────────────────────────────────────────
\section{Experiments}
\label{sec:exp}

\subsection{Experimental Setup}
\label{sec:setup}

We summarise our main experiments; all runs are available at:
\url{https://wandb.ai/fotiskapotos/rl-snake?nw=nwuserfotiskapotos}.

Our primary evaluation metric is the snake length, complemented by
total reward and episode length (steps). Due to the high variance of
training, all curves are smoothed using an exponential moving average.

\subsection{MLP agent on the easy environment}
\label{sec:exp_baseline}
We start by training an MLP agent on the easy environment.
The agent quickly learns to collect food and avoid collisions.
We then study the impact of reward design by introducing a step penalty
and distance-based shaping to accelerate learning.
While these modifications slightly improve convergence speed, the overall
performance gains remain marginal.
We also conduct a brief analysis of reward scaling by multiplying all rewards
by a constant factor. This has limited impact on the final performance,
but can affect training stability and the speed of convergence.

\begin{figure}[h]
  \centering
  \begin{subfigure}{0.48\linewidth}
    \centering
    \includegraphics[width=\linewidth]{../poster/img/MLP rewards.png}
    \caption{Baseline vs.\ step penalty vs.\ distance shaping}
  \end{subfigure}
  \hfill
  \begin{subfigure}{0.48\linewidth}
    \centering
    \includegraphics[width=\linewidth]{../poster/img/reward magnitude_scaling.png}
    \caption{Effect of reward magnitude scaling}
  \end{subfigure}
  \caption{MLP agent snake length over training (5\,000 episodes, $16\!\times\!16$ grid).}
\end{figure}

\subsection{CNN agents on the easy environment}
We next replace the MLP with convolutional agents operating directly on
grid observations. We consider both egocentric-window inputs and full-grid
representations.

\begin{figure}[h]
  \centering
  \includegraphics[width=0.7\linewidth]{../poster/img/Ergovsfull.png}
  \caption{Comparison between egocentric and full-grid CNN agents on the easy environment.}
\end{figure}

The egocentric CNN learns faster in the early stages of training, benefiting
from its focused local representation. However, it converges to a slightly
lower final performance compared to the full-grid CNN, which has access to
global spatial information.

We also train the full CNN agent for longer and compare against the performances of the MLP agent.

\begin{figure}[h]
  \centering
  \includegraphics[width=0.7\linewidth]{../poster/img/train_simple.png}
  \caption{Comparison between MLP and CNN on the easy enviroment}
\end{figure}

We find that even thought the CNN agent converges much slower, its final performance is slightly better than the MLP.

\subsection{Results on the Hard Environment}
\subsubsection{Adaptation of Pre-trained Agents}
We evaluate how agents trained in simpler settings transfer to the hard
environment. Performance drops significantly for the MLP, particularly in
the presence of dynamic obstacles, where both reward and survival degrade
and the loop rate increases sharply. This suggests that hand-crafted features
fail to capture the additional complexity and non-stationarity of the environment.

In contrast, CNN agents exhibit substantially better robustness. While their
performance also decreases under moving obstacles, the degradation remains
limited and loop behaviours are less frequent. This indicates that spatial
representations enable better adaptation to complex dynamics and obstacle
interactions.

Overall, the results highlight a clear gap in generalization: CNN-based agents
transfer more effectively to the hard regime, whereas MLP agents struggle to
cope with increased environmental complexity.

\begin{table}[h]
\centering
\small
\begin{tabular}{lcccc}
\toprule
\textbf{Agent} & \textbf{Mean length} & \textbf{Mean reward} & \textbf{Mean foods} & \textbf{Loop rate} \\
\midrule
MLP (static hard)   & 13.2 & 10.8 & 10.2 & 0.05 \\
MLP (moving obs.)   &  9.1 &  6.7 &  6.4 & 0.27 \\
CNN (static hard)   & 15.6 & 12.9 & 12.5 & 0.02 \\
CNN (moving obs.)   & 14.8 & 11.7 & 11.3 & 0.08 \\
\bottomrule
\end{tabular}
\caption{Performance of MLP and CNN agents in the hard environment.}
\label{tab:hard_results}
\end{table}
\subsubsection{Training CNN Agents on the Hard Environment}

We finally train a CNN agent directly on the hard environment.
We incorporate the full set of stabilisation techniques introduced earlier,
including Double DQN, Dueling DQN, and frame stacking, and tune the reward
parameters to improve convergence and stability.

\begin{figure}[h]
  \centering
  \includegraphics[width=0.7\linewidth]{../poster/img/finalfinal.png}
  \caption{Training performance of the CNN agent on the hard environment.}
\end{figure}


% ── 6. Conclusion ─────────────────────────────────────────────────────────────
\section{Conclusion}
\label{sec:conclusion}

We presented a comparative study of Deep Q-Network variants on a configurable Snake environment of increasing complexity. Our results show that while hand-crafted MLP agents are sufficient in simple settings, they fail to scale to more complex, dynamic environments. In contrast, CNN-based agents, particularly when combined with Double DQN, Dueling architectures, and frame stacking, demonstrate stronger performance and significantly better generalization.

Overall, our findings highlight the importance of spatial representations and stabilisation techniques for solving reinforcement learning problems with rich dynamics and partial observability.
\end{document}
